\renewcommand{\abstractname}{Resumen}

\begin{abstract}
Este Trabajo de Fin de Grado aborda las ineficiencias derivadas del proceso manual de gestión de gastos en una empresa de consultoría informática. Este proceso suponía una elevada carga administrativa y una alta tasa de errores durante la generación de los informes de gastos. Para resolver este problema, se ha diseñado y desarrollado una aplicación web full-stack que centraliza todo el flujo de trabajo en una única plataforma digital y automatiza la extracción de información de los tiquets y recibos de gasto.

Esta aplicación ha sido desarrollada utilizando .NET 8 para la capa del backend, React con TypeScript para la capa del frontend y Microsoft Azure para la provisión de recursos en la nube. La arquitectura del sistema sigue los principios de Clean Architecture y SOLID con el fin de obtener una aplicación mantenible, escalable y fácilmente integrable con servicios externos. 

Asimismo, la arquitectura combina operaciones síncronas y asíncronas. Gracias a esta estrategia, la aplicación puede ofrecer una experiencia de usuario rápida y fluida, mientras que las tareas que tienen un elevado coste computacional, como la extracción de datos de recibos mediante inteligencia artificial, se ejecutan en segundo plano.

Para garantizar el éxito del proyecto, los requisitos de la aplicación se han priorizado mediante la metodología MoSCoW, lo que ha permitido implementar los requisitos de mayor prioridad dentro del tiempo y el presupuesto establecidos. Los resultados obtenidos con la nueva aplicación demuestran que la solución resuelve eficazmente el problema identificado, obteniendo una mejora media de la eficiencia del 86,5 \%.
\end{abstract}